% from ega4-2-fr.tex
\oldpage[IV-2]{40}
\subsection{Décompositions irredondantes}

\begin{proposition}[3.2.1]
\label{IV.2.3.2.1-fr}
Soient $X$ un préschéma localement noethérien, $U$ un ouvert dense dans $X$.
Les conditions suivantes sont équivalentes~:
\begin{enumerate}
\item[\emph{a)}] $X$ est réduit.
\item[\emph{b)}] Le sous-préschéma induit sur $U$ est réduit et $X$ est sans
cycle premier immergé.
\item[\emph{c)}] $X$ est sans cycle premier immergé et pour tout point
générique $x$ d'une composante irréductible de $X$, on a
$\operatorname{long}(\mathcal O_x)=1$.
\end{enumerate}
Les cycles premiers associés à $X$ sont alors identiques aux composantes
irréductibles de $X$.
\end{proposition}

Il est clair que si $X$ est réduit il en est de même du sous-préschéma induit
sur $U$. En outre, l'existence des cycles premiers immergés étant locale, on
peut se borner à considérer le cas où $X=\operatorname{Spec}(A)$ est affine,
$A$ étant noethérien. Si $A$ est réduit, on sait que les idéaux premiers
minimaux de $A$ forment une décomposition primaire réduite de $(0)$
(Bourbaki, \emph{Alg. comm.}, chap.~IV, \S~2, n\textsuperscript{o}~5,
prop.~10) et sont les éléments de $\operatorname{Ass}(A)$, donc il n'existe
pas d'idéaux premiers immergés associés à $A$, ce qui montre que \emph{a)}
implique \emph{b)}. Il est immédiat que \emph{b)} entraîne \emph{c)}, car un
point générique $x$ d'une composante irréductible de $X$ appartient à $U$,
donc $\mathcal O_x$ est un corps. Enfin, \emph{c)} entraîne \emph{a)}~: il
suffit en effet de remarquer que si $\mathcal N$ est le Nilradical de
$\mathcal O_X$, qui est un Idéal cohérent, $\operatorname{Supp}(\mathcal N)$
ne peut contenir aucun des points génériques des composantes irréductibles de
$X$ par hypothèse~; si $\operatorname{Supp}(\mathcal N)$ n'était pas vide et
si $x$ était un des points maximaux de cet ensemble fermé, le critère
(3.1.3, \emph{c$'$)}) montrerait que $x\in\operatorname{Ass}(\mathcal O_X)$,
et $\overline{\{x\}}$ serait donc un cycle premier \emph{immergé} de $X$,
contrairement à l'hypothèse~; donc $\mathcal N=0$.

\begin{definition}[3.2.2]
\label{IV.2.3.2.2-fr}
Soient $X$ un préschéma localement noethérien, $\mathcal F$ un
$\mathcal O_X$-Module cohérent. On dit que $\mathcal F$ est \emph{réduit}
s'il vérifie les deux conditions suivantes~: 1\textsuperscript{o}
$\mathcal F$ est sans cycle premier associé immergé~; 2\textsuperscript{o}
pour tout point maximal $x$ de $\operatorname{Supp}(\mathcal F)$, on a
$\operatorname{long}(\mathcal F_x)=1$.
\end{definition}

La condition 1\textsuperscript{o} signifie que les cycles premiers associés à
$\mathcal F$ sont les composantes irréductibles de
$\operatorname{Supp}(\mathcal F)$ (3.1.4), et la condition
2\textsuperscript{o} signifie que pour tout point générique $x$ d'une telle
composante on a $\operatorname{long}(\mathcal F_x)=1$.

Pour un schéma affine $X$, cette définition donne en particulier la notion de
\emph{module réduit} sur un anneau noethérien $A$~; un $A$-module de type fini
$M$ est dit \emph{réduit}, s'il n'a pas d'idéaux premiers associés immergés et
si, pour tout $\mathfrak p\in\operatorname{Ass}(M)$,
$\operatorname{long}_{A_\mathfrak p}(M_\mathfrak p)=1$. Revenant au cas d'un
préschéma localement noethérien $X$ et d'un $\mathcal O_X$-Module cohérent
$\mathcal F$, on dit que $\mathcal F$ est \emph{réduit en un point} $x\in X$
si $\mathcal F_x$ est un $\mathcal O_x$-module réduit~; cela signifie encore
que, sur le schéma local $\operatorname{Spec}(\mathcal O_x)$,
$\widetilde{\mathcal F_x}$ est réduit~; il revient donc au même de dire
\emph{que $x$ n'appartient à aucun cycle premier associé immergé de
$\mathcal F$ et que $\operatorname{long}(\mathcal F_z)=1$ pour tous les
points maximaux $z$ de $\operatorname{Supp}(\mathcal F)$ tels que
$x\in\overline{\{z\}}$}. Il est clair que si $\mathcal F$ est un
$\mathcal O_X$-Module cohérent réduit, il est réduit en chaque point de $X$~;
inversement, si $\mathcal F$ est \emph{réduit au point} $x$, il existe un
voisinage ouvert $U$ de $x$ tel que $\mathcal F|U$ soit un
$\mathcal O_U$-Module \emph{réduit}~: il suffit en effet de prendre $U$ ne
rencontrant aucun cycle premier associé immergé de $\mathcal F$ (un tel
voisinage existe puisque ces cycles forment un ensemble localement fini de
parties
